%% ============================================================
%%  S+SSPR 2026 talk — Orthogonality and Dimensionality in
%%  Airline Cluster Analysis using PCA and Kernel PCA
%%  Story line: "The linearity certificate"
%%  Compile: lualatex talk.tex   (see Makefile)
%% ============================================================
\documentclass[10pt,aspectratio=169]{beamer}

\usepackage{asc-theme}
\usepackage{amsmath,amssymb}
\usepackage{booktabs}
\usepackage{pgfplots}
\pgfplotsset{compat=1.18}
\usetikzlibrary{calc, positioning, arrows.meta}

%% ── Helper commands (as in the paper) ───────────────────────
\newcommand{\pc}[1]{PC#1}
\newcommand{\cl}[1]{$C_{#1}$}
\newcommand{\kmeans}{$k$-means}
\newcommand{\ari}{ARI}
\newcommand{\pr}{\textit{PR}}
\newcommand{\sil}{\textit{Sil}}
\newcommand{\db}{\textit{DB}}
\newcommand{\kpca}{KPCA}

%% ── Plot chrome in asc palette; cluster colors as in paper ──
\pgfplotsset{strip/.style={
  width=1.08\textwidth, height=3.1cm,
  ytick={1,2,3,4,5,6},
  yticklabels={$C_5$,$C_4$,$C_3$,$C_2$,$C_1$,$C_0$},
  ymin=0.5, ymax=6.5,
  tick label style={font=\fontsize{4.5}{5}\selectfont, color=ascSlate600},
  title style={font=\tiny\bfseries, color=ascSlate900},
}}
\pgfplotsset{asc/.style={
  axis line style={ascSlate400},
  every tick/.append style={ascSlate400},
  tick label style={font=\tiny, color=ascSlate600},
  label style={font=\scriptsize, color=ascSlate600},
  title style={font=\scriptsize\bfseries, color=ascSlate900},
  grid=both,
  grid style={line width=0.25pt, draw=ascSlate100},
  legend style={font=\tiny, draw=ascSlate100, fill=white, text=ascSlate600,
                inner sep=2.5pt, row sep=0.5pt},
}}

\title[Orthogonality \& Dimensionality in Airline Cluster Analysis]%
  {Orthogonality and Dimensionality in Airline Cluster Analysis using PCA and Kernel PCA}
\author{Andreas Schlapbach}
\date{S+SSPR 2026 --- Bern, August 24--26}

\begin{document}


\begin{frame}[plain]
  \vspace*{0.9cm}
  {\plexmedium\small\color{ascTeal600}S+SSPR 2026 --- Bern, August 24--26\par}
  \vspace{0.35cm}
  {\usebeamerfont{title}\usebeamercolor[fg]{title}%
    Orthogonality and Dimensionality in Airline Cluster Analysis using PCA and Kernel PCA\par}
  \vspace{0.18cm}
  {\color{ascTeal500}\rule{28pt}{2.4pt}\par}
  \vspace{0.30cm}
  {\large\muted{The linearity certificate: how to show a dataset needs no kernel}\par}
  \vspace{0.85cm}
  {\usebeamerfont{author}\usebeamercolor[fg]{author}Andreas Schlapbach\par}
  \vspace{0.10cm}
  {\usebeamerfont{institute}\usebeamercolor[fg]{institute}%
    Swiss Federal Railways (SBB), Switzerland\footnotemark[1]\quad\textbullet\quad
    {\color{ascTeal600}asc.engineering}
    }
  \vspace{0.55cm}
  \footnotetext[1]{\tiny\color{ascSlate400}The views expressed are the author's
    own and do not necessarily reflect the views or policies of SBB.}
\end{frame}



\begin{frame}{Kernel PCA turns 28}
  \begin{itemize}
     \item{The kernel methods turn 28 age this year:  \muted{Sch\"olkopf, Smola \& M\"uller, \emph{Nonlinear component analysis as a kernel eigenvalue problem}, Neural Computation, 1998.}}

    \item {But {\color{ascTeal600}How do you certify that a dataset contains none?}}

    \item {A certified \emph{no} is a result: it licenses every linear tool
	    downstream --- PCA, $k$-means in score space, linear ridge.}
    \item {Papers rarely issue that certificate explicitly. This talk shows
	    how, on a published study.}
  \end{itemize}
\end{frame}


\begin{frame}{The Case Under Audit}
   \begin{itemize}
      \item \emph{Renold et al.: "Methodological Framework for a Deeper Understanding of Airline Proift Cycles in the Context of Disruptive Exogenous Impact", Journal of Air Transport Management, 2023}
      \item Studies US airline profit cycles between 1995 and 2020
      \item Data set consists of seven variables: operating profit, fuel, wages, RPM,
          load factor, yield, other expenses
      \item \kmeans{} with $k=6$, PCA, system-dynamics modelling
      \item \emph{Headline claim:} COVID-19 does \emph{not} permanently 
	  disrupt airline profit regimes: post-pandemic patterns revert
          to pre-COVID clusters.
    \end{itemize}
\end{frame}


\begin{frame}{Two Risks Left Unexamined}
   \begin{enumerate}
      \item The seven inputs are \teal{highly collinear} ---
        secular growth drives them all. Does that distort the Euclidean
        distances \kmeans{} partitions by?
      \item The choice \teal{$k=6$ has no quantitative justification} ---
          no quality criterion, no alternative $k$ compared.
    \end{enumerate}
    \vspace{0.15cm}
     \muted{If either is material, the regime labels --- and the headline
      claim --- cannot be taken at face value.}
\end{frame}


\begin{frame}{Seven Variables, 2.46 Effective Dimensions}
  The first measurement: the \textbf{participation ratio} of the eigenvalue
  spectrum to see how many dimensions actually carry signal.
  \vspace{0.2cm}
  \begin{equation*}
    \pr \;=\; \frac{\bigl(\sum_i \lambda_i\bigr)^2}{\sum_i \lambda_i^2}
        \;=\; \frac{100^2}{59.8^2 + 16.7^2 + 11.3^2 + 7.8^2 + 3.9^2 + 0.4^2 + 0.1^2}
        \;\approx\; {\color{ascTeal600}\mathbf{2.46}}
  \end{equation*}
  \vspace{0.25cm}
  \begin{itemize}
    \item Of seven nominal dimensions, $\sim$2.5 are real: \teal{severe
      collinearity}, one dominant shared growth factor.
    \item \textbf{Question:} Does a low \pr{} \emph{predict} a near-linear, low-dimensional manifold
    \item \textbf{Plan:} establish the linear geometry, then attack it
      with every kernel family we have. If the linear story survives, it is
      certified.
  \end{itemize}
\end{frame}

\begin{frame}{The Three Spaces}
  All 26 standardized annual observations, clustered in three representations:
  
  \begin{center}
  \small
  \begin{tabular}{@{}clrr@{}}
  \toprule
  \textbf{\pc{}} & \textbf{Interpretation} &
  \textbf{Variance (\%)} & \textbf{Cumulative (\%)} \\
  \midrule
  1 & Secular US aviation growth     & 59.8 & 59.8 \\
  2 & Profit/wage boom--bust tension & 16.7 & 76.5 \\
  3 & Fuel-price shocks              & 11.3 & 87.8 \\
  4 & \teal{COVID-era yield collapse} &  7.8 & 95.6 \\
  5--7 & Noise                       &  4.4 & 100.0 \\
  \bottomrule
  \end{tabular}
  \end{center}
  \vspace{0.35cm}
  \begin{columns}[T]
    \begin{column}{0.31\textwidth}
      \centering\textbf{7D-raw}\\[2pt]
      \muted{\small original variables,\\collinearity included}
    \end{column}
    \begin{column}{0.31\textwidth}
      \centering\textbf{3D-\pc{}}\\[2pt]
      \muted{\small 87.8\% of variance,\\geometrically clean}
    \end{column}
    \begin{column}{0.31\textwidth}
      \centering\textbf{4D-\pc{}}\\[2pt]
      \muted{\small 95.6\%, adds the\\COVID yield dimension}
    \end{column}
  \end{columns}
\end{frame}

\begin{frame}{Finding I: Assignments Are Invariant Across Spaces}
  \begin{columns}[T]
    \begin{column}{0.56\textwidth}
      \scriptsize
      \setlength{\tabcolsep}{4pt}
      \begin{tabular}{@{}l l c c c@{}}
      \toprule
      \textbf{Years} & \textbf{Regime} &
      \textbf{7D-raw} & \textbf{3D-\pc{}} & \textbf{4D-\pc{}} \\
      \midrule
      1995--1999 & 9/11 era (early) & \cl{2} & \cl{2} & \cl{2} \\
      2000--2005 & 9/11 + post-9/11 & \cl{1} & \cl{1} & \cl{1}/\cl{4} \\
      2006--2010 & Financial Crisis & \cl{4} & \cl{4} & \cl{4} \\
      2011--2014 & Post-FC          & \cl{5} & \cl{5} & \cl{5} \\
      2015--2019 & Pre-/post-COVID  & \cl{0} & \cl{0} & \cl{0} \\
      2020       & COVID            & \cl{3} & \cl{3} & \cl{3} \\
      \midrule
      \multicolumn{2}{@{}l}{\textit{Agreement vs.\ 7D-raw}} &
        --- & 26/26 & 25/26 \\
      \multicolumn{2}{@{}l}{\textit{\ari{} vs.\ 7D-raw}} &
        --- & \textbf{1.000} & 0.884 \\
      \bottomrule
      \end{tabular}
      \vspace{0.25cm}

      \small $k=6$ in 3D-\pc{} space reproduces the 7D-raw partition
      \teal{bit for bit}.
    \end{column}
    \begin{column}{0.40\textwidth}
      \textbf{Why doesn't collinearity rearrange anything?}
      \vspace{0.1cm}
      \begin{itemize}
        \item It stretches distances $\sim$3$\times$ \emph{along} \pc{1} ---
          without reordering points \emph{within} it.
        \item The six clusters are contiguous \emph{time bands} along \pc{1}:
          stretched, but no year crosses a boundary.
        \item Cross-sectional clusters (legacy vs.\ low-cost carriers) would
          cut \emph{across} the growth axis --- there, distortion would be
          severe.
      \end{itemize}
    \end{column}
  \end{columns}
\end{frame}

\begin{frame}{The Geometry Behind the Invariance}
  \begin{columns}[T]
    \begin{column}{0.60\textwidth}
      \centering
      \begin{tikzpicture}
      \begin{axis}[asc,
        view={120}{30},
        width=\textwidth, height=5.6cm,
        xlabel=\pc{1} (59.8\%), ylabel=\pc{2} (16.7\%), zlabel=\pc{3} (11.3\%),
        xmin=-3, xmax=3, ymin=-2.5, ymax=2.5, zmin=-2, zmax=2,
        every axis plot/.append style={line width=1.2pt, mark size=2.2pt},
      ]
      \addplot3[color=blue, mark=diamond*, only marks] coordinates {
        (-2.8,-0.8,-1.2) (-2.5,-0.5,-1.0) (-2.7,-0.7,-1.1)
        (-2.9,-0.9,-1.3) (-2.6,-0.6,-0.9)};
      \addplot3[color=orange, mark=square*, only marks] coordinates {
        (-1.5,1.8,0.4) (-1.3,1.9,0.3) (-1.6,1.7,0.5)
        (-1.2,1.6,0.2) (-1.4,1.8,0.4) (-1.5,1.9,0.3)};
      \addplot3[color=green!70!black, mark=pentagon*, only marks] coordinates {
        (0.8,0.9,-0.6) (0.5,1.0,-0.5) (0.9,0.8,-0.7)
        (1.1,1.1,-0.8) (0.6,0.7,-0.4)};
      \addplot3[color=brown, mark=asterisk, only marks] coordinates {
        (1.6,0.2,-0.3) (1.7,0.1,-0.2) (1.5,0.0,-0.4) (1.8,0.3,-0.1)};
      \addplot3[color=red, mark=*, only marks] coordinates {
        (2.1,0.5,-0.2) (2.3,0.4,0.1) (2.2,0.6,0.0)
        (2.0,0.3,-0.1) (2.4,0.7,0.2)};
      \addplot3[color=purple, mark=triangle*, only marks] coordinates {
        (2.2,-1.8,1.6)};
      \end{axis}
      \end{tikzpicture}

      {\scriptsize\muted{3D-\pc{} space, 87.8\% of variance. Chronological
      bands left to right:
      {\color{blue}1995--99}, {\color{orange}2000--05},
      {\color{green!70!black}2006--10}, {\color{brown}2011--14},
      {\color{red}2015--19}; {\color{purple}2020} displaced.}}
    \end{column}
    \begin{column}{0.38\textwidth}
      \vspace{0.3cm}
      \begin{itemize}
        \item The six-cluster taxonomy is a \teal{temporal partition} along
          \pc{1} --- early losses $\to$ recovery $\to$ peak profit.
        \item Secondary structure: profit--wage tension (\pc{2}), fuel shocks
          (\pc{3}).
        \item COVID 2020 is a \teal{geometric singularity}, displaced by the
          yield collapse (load factor $85\% \to 59\%$).
      \end{itemize}
      \vspace{0.2cm}
      \muted{\small This is the linear story. Now we stress-test it.}
    \end{column}
  \end{columns}
\end{frame}

\begin{frame}{Finding II: The Clean Space Reveals $k=3$}
  \begin{center}
  \begin{tikzpicture}
  \begin{axis}[asc,
    name=silplot,
    width=0.46\textwidth, height=4.4cm,
    xlabel={$k$ (number of clusters)}, ylabel={Silhouette coefficient},
    xmin=1.5, xmax=9.5, ymin=0.38, ymax=0.56,
    xtick={2,3,4,5,6,7,8,9},
    legend style={at={(0.98,0.05)}, anchor=south east},
  ]
  \addplot[red!80!black, thick, solid, mark=*, mark size=1.8pt]
    coordinates {(2,0.431)(3,0.443)(4,0.429)(5,0.413)(6,0.448)(7,0.449)(8,0.439)(9,0.416)};
  \addlegendentry{7D-raw}
  \addplot[blue!80!black, thick, dashed, mark=square*, mark size=1.8pt]
    coordinates {(2,0.506)(3,0.523)(4,0.481)(5,0.421)(6,0.448)(7,0.429)(8,0.416)(9,0.430)};
  \addlegendentry{3D-\pc{}}
  \addplot[ascSlate400, thick, dotted, mark=triangle*, mark size=1.8pt]
    coordinates {(2,0.458)(3,0.472)(4,0.466)(5,0.460)(6,0.490)(7,0.463)(8,0.449)(9,0.410)};
  \addlegendentry{4D-\pc{}}
  \draw[ascTeal500, very thick] (axis cs:3,0.38) -- (axis cs:3,0.523);
  \node[font=\tiny\bfseries, anchor=south, color=ascTeal600] at (axis cs:3,0.523) {$k^*\!=\!3$};
  \end{axis}
  \begin{axis}[asc,
    at={(silplot.south east)}, xshift=1.3cm, anchor=south west,
    width=0.46\textwidth, height=4.4cm,
    xlabel={$k$ (number of clusters)}, ylabel={Davies--Bouldin index},
    xmin=1.5, xmax=9.5, ymin=0.45, ymax=0.96,
    xtick={2,3,4,5,6,7,8,9},
    legend style={at={(0.98,0.98)}, anchor=north east},
  ]
  \addplot[red!80!black, thick, solid, mark=*, mark size=1.8pt]
    coordinates {(2,0.906)(3,0.648)(4,0.658)(5,0.679)(6,0.647)(7,0.629)(8,0.613)(9,0.538)};
  \addlegendentry{7D-raw}
  \addplot[blue!80!black, thick, dashed, mark=square*, mark size=1.8pt]
    coordinates {(2,0.742)(3,0.507)(4,0.603)(5,0.603)(6,0.584)(7,0.583)(8,0.596)(9,0.554)};
  \addlegendentry{3D-\pc{}}
  \addplot[ascSlate400, thick, dotted, mark=triangle*, mark size=1.8pt]
    coordinates {(2,0.853)(3,0.601)(4,0.615)(5,0.569)(6,0.557)(7,0.640)(8,0.627)(9,0.578)};
  \addlegendentry{4D-\pc{}}
  \draw[ascTeal500, very thick] (axis cs:3,0.45) -- (axis cs:3,0.507);
  \node[font=\tiny\bfseries, anchor=north, color=ascTeal600] at (axis cs:3,0.507) {$k^*\!=\!3$};
  \end{axis}
  \end{tikzpicture}
  \end{center}
  \vspace{0.1cm}
  \begin{itemize}
    \item In 3D-\pc{} space the silhouette peaks at \teal{$k=3$}, $k=6$ scores 14\% worse. In 7D-raw the curve is
      \teal{flat}: the inflated growth axis suppresses the contrast.
    \item Three regimes: low-profit 1995--2005, recovery 2006--14, peak
      2015--19 (+ COVID outlier). $k=6$ is coherent but \emph{imposed}.
    \item \muted{Suggestive, not conclusive: two internal indices, $n=26$, no
      significance test.}
  \end{itemize}
\end{frame}


\begin{frame}{Enter the Kernels: the Certificate's Core Idea}
  \textbf{If the manifold hides curvature, \emph{some} kernel family should
  exploit it.} So we attack with all of them --- on the 4D-\pc{} scores,
  assignments held fixed.
  \vspace{0.35cm}
  \begin{columns}[T]
    \begin{column}{0.31\textwidth}
      \begin{exampleblock}{Distance-based}
        \small RBF ($\gamma$ by median pairwise distance)\\[2pt]
        Polynomial, degree 2\\[2pt]
        Polynomial, degree 3
      \end{exampleblock}
    \end{column}
    \begin{column}{0.31\textwidth}
      \begin{exampleblock}{Information-geometric}
        \small Fisher--Mahalanobis:\\
        $K_F(x,y) = x^{\top}\hat\Sigma^{-1}y$\\[2pt]
        = PCA on whitened scores; amplifies \pc{4} by
        $1/0.078 \approx 12.8\times$
      \end{exampleblock}
    \end{column}
    \begin{column}{0.31\textwidth}
      \begin{exampleblock}{Graph-based}
        \small Diffusion $K=\exp(-\beta L)$\\
        on the 5-NN graph\\[2pt]
        local connectivity instead of a distance formula
      \end{exampleblock}
    \end{column}
  \end{columns}
  \vspace{0.25cm}
  \begin{itemize}
    \item Plus the \textbf{linear kernel} as baseline ($\equiv$ standard PCA).
    \item All bandwidths by \teal{median heuristics} --- nothing fitted,
      nothing tuned toward a verdict.
    \item Three families fail for three \emph{different} reasons --- their
      agreement is evidence.
  \end{itemize}
\end{frame}



\begin{frame}{Linear PCA vs.\ RBF kernel PCA}
  \begin{columns}[T]
    \begin{column}{0.48\textwidth}
      \centering
      \begin{tikzpicture}
      \begin{axis}[asc,
        width=\textwidth, height=4.9cm,
        xlabel=\pc{1} (59.8\%), ylabel=\pc{2} (16.7\%),
        xmin=-3.2, xmax=2.8, ymin=-2.2, ymax=2.2,
        title={Linear PCA},
      ]
      \addplot[ascSlate400, dashed, line width=0.5pt] coordinates {(-1.9,-2.2) (-1.9,2.2)};
      \addplot[ascSlate400, dashed, line width=0.5pt] coordinates {(-0.2,-2.2) (-0.2,2.2)};
      \addplot[ascSlate400, dashed, line width=0.5pt] coordinates {(1.2,-2.2) (1.2,2.2)};
      \addplot[ascSlate400, dashed, line width=0.5pt] coordinates {(1.85,-2.2) (1.85,2.2)};
      \addplot[color=red, mark=*, only marks, mark size=2pt] coordinates {
        (2.1,0.5) (2.3,0.4) (2.2,0.6) (2.0,0.3) (2.4,0.7)};
      \addplot[color=orange, mark=square*, only marks, mark size=2pt] coordinates {
        (-1.5,1.8) (-1.3,1.9) (-1.6,1.7) (-1.2,1.6) (-1.4,1.8) (-1.5,1.9)};
      \addplot[color=blue, mark=diamond*, only marks, mark size=2pt] coordinates {
        (-2.8,-0.8) (-2.5,-0.5) (-2.7,-0.7) (-2.9,-0.9) (-2.6,-0.6)};
      \addplot[color=purple, mark=triangle*, only marks, mark size=2.4pt] coordinates {
        (2.2,-1.8)};
      \addplot[color=green!70!black, mark=pentagon*, only marks, mark size=2pt] coordinates {
        (0.8,0.9) (0.5,1.0) (0.9,0.8) (1.1,1.1) (0.6,0.7)};
      \addplot[color=brown, mark=asterisk, only marks, mark size=2pt] coordinates {
        (1.6,0.2) (1.7,0.1) (1.5,0.0) (1.8,0.3)};
      \end{axis}
      \end{tikzpicture}
    \end{column}
    \begin{column}{0.48\textwidth}
      \centering
      \begin{tikzpicture}
      \begin{axis}[asc,
        width=\textwidth, height=4.9cm,
        xlabel=\kpca{}1, ylabel=\kpca{}2,
        xmin=-2.0, xmax=1.8, ymin=-1.7, ymax=1.5,
        title={RBF kernel PCA ($\gamma = \gamma_{\mathrm{med}}$)},
      ]
      \addplot[ascSlate400, dashed, line width=0.5pt] coordinates {(-1.25,-1.7) (-1.25,1.5)};
      \addplot[ascSlate400, dashed, line width=0.5pt] coordinates {(-0.35,-1.7) (-0.35,1.5)};
      \addplot[ascSlate400, dashed, line width=0.5pt] coordinates {(0.25,-1.7) (0.25,1.5)};
      \addplot[ascSlate400, dashed, line width=0.5pt] coordinates {(0.8,-1.7) (0.8,1.5)};
      \addplot[color=red, mark=*, only marks, mark size=2pt] coordinates {
        (1.2,0.05) (1.3,0.1) (1.1,0.15) (1.4,0.0) (1.5,0.25)};
      \addplot[color=orange, mark=square*, only marks, mark size=2pt] coordinates {
        (-0.95,0.9) (-0.85,1.0) (-1.1,0.8) (-0.8,0.95) (-1.0,0.85) (-0.9,1.05)};
      \addplot[color=blue, mark=diamond*, only marks, mark size=2pt] coordinates {
        (-1.65,-0.55) (-1.5,-0.4) (-1.7,-0.5) (-1.8,-0.65) (-1.6,-0.45)};
      \addplot[color=purple, mark=triangle*, only marks, mark size=2.4pt] coordinates {
        (0.6,-1.6)};
      \addplot[color=green!70!black, mark=pentagon*, only marks, mark size=2pt] coordinates {
        (0.15,0.4) (0.05,0.5) (0.25,0.35) (0.35,0.55) (0.1,0.45)};
      \addplot[color=brown, mark=asterisk, only marks, mark size=2pt] coordinates {
        (0.75,0.0) (0.8,-0.05) (0.65,0.05) (0.9,-0.1)};
      \end{axis}
      \end{tikzpicture}
    \end{column}
  \end{columns}
  \vspace{0.15cm}
  \begin{itemize}
    \item \teal{\ari{} = 1.0}: all 26 observations land in identical
      clusters under the kernel map.
    \item COVID ({\color{purple}$\blacktriangle$}) becomes \emph{more}
      isolated --- \pc{4}'s yield collapse enters the pairwise distances of
      the Gram matrix. $\pr \approx 2.46$ predicted this; \kpca{} confirms it.
  \end{itemize}
\end{frame}


\begin{frame}{Six Kernels, One Verdict: the \kpca{}1 Strip Charts}
  \begin{columns}[T]
  \begin{column}{0.32\textwidth}\centering
    \begin{tikzpicture}\begin{axis}[asc, strip, title={Linear (= PCA baseline)}, xmin=-3.5, xmax=3.0]
    \addplot[fill=ascSlate100, draw=none] coordinates {(1.95,2.5)(2.45,2.5)(2.45,6.5)(1.95,6.5)} -- cycle;
    \addplot[color=red,mark=*,only marks,mark size=1.4pt] coordinates {(2.1,6)(2.3,6)(2.2,6)(2.0,6)(2.4,6)};
    \addplot[color=orange,mark=square*,only marks,mark size=1.4pt] coordinates {(-1.5,5)(-1.3,5)(-1.6,5)(-1.2,5)(-1.4,5)(-1.5,5)};
    \addplot[color=blue,mark=diamond*,only marks,mark size=1.4pt] coordinates {(-2.8,4)(-2.5,4)(-2.7,4)(-2.9,4)(-2.6,4)};
    \addplot[color=purple,mark=triangle*,only marks,mark size=1.8pt] coordinates {(2.2,3)};
    \addplot[color=green!70!black,mark=pentagon*,only marks,mark size=1.4pt] coordinates {(0.8,2)(0.5,2)(0.9,2)(1.1,2)(0.6,2)};
    \addplot[color=brown,mark=asterisk,only marks,mark size=1.4pt] coordinates {(1.6,1)(1.7,1)(1.5,1)(1.8,1)};
    \end{axis}\end{tikzpicture}
  \end{column}
  \begin{column}{0.32\textwidth}\centering
    \begin{tikzpicture}\begin{axis}[asc, strip, title={RBF}, xmin=-2.2, xmax=2.1]
    \addplot[fill=ascSlate100, draw=none] coordinates {(0.57,0.5)(0.93,0.5)(0.93,3.5)(0.57,3.5)} -- cycle;
    \addplot[color=red,mark=*,only marks,mark size=1.4pt] coordinates {(1.2,6)(1.3,6)(1.1,6)(1.4,6)(1.5,6)};
    \addplot[color=orange,mark=square*,only marks,mark size=1.4pt] coordinates {(-0.95,5)(-0.85,5)(-1.10,5)(-0.80,5)(-1.00,5)(-0.90,5)};
    \addplot[color=blue,mark=diamond*,only marks,mark size=1.4pt] coordinates {(-1.65,4)(-1.50,4)(-1.70,4)(-1.80,4)(-1.60,4)};
    \addplot[color=purple,mark=triangle*,only marks,mark size=1.8pt] coordinates {(0.60,3)};
    \addplot[color=green!70!black,mark=pentagon*,only marks,mark size=1.4pt] coordinates {(0.15,2)(0.05,2)(0.25,2)(0.35,2)(0.10,2)};
    \addplot[color=brown,mark=asterisk,only marks,mark size=1.4pt] coordinates {(0.75,1)(0.80,1)(0.65,1)(0.90,1)};
    \end{axis}\end{tikzpicture}
  \end{column}
  \begin{column}{0.32\textwidth}\centering
    \begin{tikzpicture}\begin{axis}[asc, strip, title={Polynomial (deg.\ 2)}, xmin=-2.0, xmax=1.8]
    \addplot[fill=ascSlate100, draw=none] coordinates {(0.52,0.5)(0.83,0.5)(0.83,3.5)(0.52,3.5)} -- cycle;
    \addplot[color=red,mark=*,only marks,mark size=1.4pt] coordinates {(0.95,6)(1.05,6)(0.90,6)(1.15,6)(1.20,6)};
    \addplot[color=orange,mark=square*,only marks,mark size=1.4pt] coordinates {(-0.80,5)(-0.70,5)(-0.95,5)(-0.65,5)(-0.85,5)(-0.75,5)};
    \addplot[color=blue,mark=diamond*,only marks,mark size=1.4pt] coordinates {(-1.45,4)(-1.30,4)(-1.50,4)(-1.60,4)(-1.40,4)};
    \addplot[color=purple,mark=triangle*,only marks,mark size=1.8pt] coordinates {(0.55,3)};
    \addplot[color=green!70!black,mark=pentagon*,only marks,mark size=1.4pt] coordinates {(0.12,2)(0.04,2)(0.22,2)(0.28,2)(0.08,2)};
    \addplot[color=brown,mark=asterisk,only marks,mark size=1.4pt] coordinates {(0.68,1)(0.72,1)(0.60,1)(0.80,1)};
    \end{axis}\end{tikzpicture}
  \end{column}
  \end{columns}
  \vspace{0.12cm}
  \begin{columns}[T]
  \begin{column}{0.32\textwidth}\centering
    \begin{tikzpicture}\begin{axis}[asc, strip, title={Polynomial (deg.\ 3)}, xmin=-2.1, xmax=2.0]
    \addplot[fill=ascSlate100, draw=none] coordinates {(0.55,0.5)(0.88,0.5)(0.88,3.5)(0.55,3.5)} -- cycle;
    \addplot[color=red,mark=*,only marks,mark size=1.4pt] coordinates {(1.05,6)(1.15,6)(1.00,6)(1.30,6)(1.40,6)};
    \addplot[color=orange,mark=square*,only marks,mark size=1.4pt] coordinates {(-0.88,5)(-0.77,5)(-1.00,5)(-0.72,5)(-0.92,5)(-0.82,5)};
    \addplot[color=blue,mark=diamond*,only marks,mark size=1.4pt] coordinates {(-1.55,4)(-1.40,4)(-1.60,4)(-1.70,4)(-1.50,4)};
    \addplot[color=purple,mark=triangle*,only marks,mark size=1.8pt] coordinates {(0.58,3)};
    \addplot[color=green!70!black,mark=pentagon*,only marks,mark size=1.4pt] coordinates {(0.10,2)(0.02,2)(0.20,2)(0.28,2)(0.06,2)};
    \addplot[color=brown,mark=asterisk,only marks,mark size=1.4pt] coordinates {(0.70,1)(0.76,1)(0.62,1)(0.85,1)};
    \end{axis}\end{tikzpicture}
  \end{column}
  \begin{column}{0.32\textwidth}\centering
    \begin{tikzpicture}\begin{axis}[asc, strip, title={Fisher}, xmin=-2.0, xmax=1.8]
    \addplot[fill=ascSlate100, draw=none] coordinates {(0.47,0.5)(0.78,0.5)(0.78,3.5)(0.47,3.5)} -- cycle;
    \addplot[color=red,mark=*,only marks,mark size=1.4pt] coordinates {(0.90,6)(1.00,6)(0.85,6)(1.15,6)(1.20,6)};
    \addplot[color=orange,mark=square*,only marks,mark size=1.4pt] coordinates {(-0.70,5)(-0.60,5)(-0.85,5)(-0.55,5)(-0.75,5)(-0.65,5)};
    \addplot[color=blue,mark=diamond*,only marks,mark size=1.4pt] coordinates {(-1.35,4)(-1.20,4)(-1.40,4)(-1.50,4)(-1.30,4)};
    \addplot[color=purple,mark=triangle*,only marks,mark size=1.8pt] coordinates {(0.50,3)};
    \addplot[color=green!70!black,mark=pentagon*,only marks,mark size=1.4pt] coordinates {(0.12,2)(0.05,2)(0.22,2)(0.28,2)(0.08,2)};
    \addplot[color=brown,mark=asterisk,only marks,mark size=1.4pt] coordinates {(0.62,1)(0.68,1)(0.55,1)(0.75,1)};
    \end{axis}\end{tikzpicture}
  \end{column}
  \begin{column}{0.32\textwidth}\centering
    \begin{tikzpicture}\begin{axis}[asc, strip, title={Graph diffusion}, xmin=-1.5, xmax=1.3]
    \addplot[fill=ascSlate100, draw=none] coordinates {(0.32,0.5)(0.58,0.5)(0.58,3.5)(0.32,3.5)} -- cycle;
    \addplot[color=red,mark=*,only marks,mark size=1.4pt] coordinates {(0.75,6)(0.82,6)(0.70,6)(0.90,6)(0.95,6)};
    \addplot[color=orange,mark=square*,only marks,mark size=1.4pt] coordinates {(-0.52,5)(-0.42,5)(-0.65,5)(-0.38,5)(-0.58,5)(-0.48,5)};
    \addplot[color=blue,mark=diamond*,only marks,mark size=1.4pt] coordinates {(-0.98,4)(-0.85,4)(-1.02,4)(-1.10,4)(-0.92,4)};
    \addplot[color=purple,mark=triangle*,only marks,mark size=1.8pt] coordinates {(0.40,3)};
    \addplot[color=green!70!black,mark=pentagon*,only marks,mark size=1.4pt] coordinates {(0.10,2)(0.05,2)(0.18,2)(0.22,2)(0.08,2)};
    \addplot[color=brown,mark=asterisk,only marks,mark size=1.4pt] coordinates {(0.42,1)(0.48,1)(0.35,1)(0.55,1)};
    \end{axis}\end{tikzpicture}
  \end{column}
  \end{columns}
  \vspace{0.1cm}
  {\scriptsize\muted{Shaded band: the overlap where \kmeans{} on \kpca{}1
  alone cannot separate the pair. Linear kernel: COVID
  {\color{purple}$C_3$} sits inside the {\color{red}$C_0$} band. All five
  nonlinear kernels move $C_3$ to the same {\color{brown}$C_5$} overlap.}}
\end{frame}



\begin{frame}{Reading the Agreement}
  \textbf{The same shift, three different mechanisms:}
  \vspace{0.2cm}
  \begin{itemize}
    \item \textbf{RBF / polynomial} --- curvature fitting in feature space.
    \item \textbf{Fisher} --- eigenvalue-inverse weighting: \pc{4} amplified
      $12.8\times$, enlarging COVID's yield-collapse signal.
    \item \textbf{Graph diffusion} --- pure 5-NN connectivity: similarity mass
      routed to $C_5$, the actual graph neighbors, overriding the global
      \pc{1} alignment with $C_0$.
  \end{itemize}
  \vspace{0.35cm}
  \begin{exampleblock}{The certificate}
    Six kernels from three fundamentally different families agree with each
    other --- and with linear PCA (\ari{} = 1.0 in 2D; \kpca{}2 resolves the
    $C_3$--$C_5$ ambiguity exactly as \pc{2} does). \textbf{The manifold
    contains no detectable curvature.} The absence of nonlinear structure is
    itself the finding.
  \end{exampleblock}
\end{frame}


\begin{frame}{The Supervised Complement: Kernel Ridge Regression}
  Structure said linear. Does nonlinear \emph{prediction} disagree?
  Predict operating profit from the six other raw variables, LOOCV,
  nested grid search for $\gamma, \alpha$.
  \vspace{0.3cm}
  \begin{center}
  \small
  \begin{tabular}{@{}l r r@{}}
  \toprule
  \textbf{Model} & \textbf{$R^2$ (all years)} & \textbf{$R^2$ (excl.\ 2020)} \\
  \midrule
  Linear ridge                  & \textbf{0.881} & 0.831 \\
  Kernel ridge (RBF)            & 0.592 & \textbf{0.836} \\
  Kernel ridge (poly., deg.~2)  & 0.380 & 0.785 \\
  \bottomrule
  \end{tabular}
  \end{center}
  \vspace{0.3cm}
  \begin{itemize}
    \item \textbf{2020 excluded:} all three models comparable
      ($R^2 = 0.79$--$0.84$) --- kernel flexibility buys
      \teal{no predictive advantage}. In-distribution, the manifold is
      linear.
    \item \textbf{2020 included:} kernels collapse, linear ridge does not.
      2020's 65\% RPM/load-factor collapse sits far outside the training
      manifold; similarity-based prediction reverts to the mean.
    \item \muted{That collapse is kernel \emph{extrapolation} failing at a
      genuine outlier --- not evidence of hidden curvature.}
  \end{itemize}
\end{frame}



\begin{frame}{The Verdict on the Audit}
  \begin{itemize}
    \item[{\color{ascTeal600}\checkmark}] \textbf{The headline claim survives.}
      Cluster assignments are invariant across spaces (\ari{} = 1.0);
      the geometry behind Renold et al.'s regime labels is sound ---
      collinearity's growth axis \emph{aligns} with the temporal partition
      instead of cutting across it.
    \item[{\color{ascTeal600}\checkmark}] \textbf{The manifold is certifiably
      linear.} Six kernels, three families, one verdict --- reinforced by
      kernel ridge showing no nonlinear accuracy gain in-distribution.
    \item[{\color{ascSlate400}$\circ$}] \textbf{But the taxonomy is
      over-resolved.} The data structurally supports \teal{three} regimes;
      $k=6$ is imposed, defensible for year-by-year policy description, not
      recommended by the indices.
    \item[{\color{ascSlate400}$\circ$}] \textbf{Practical recommendation:}
      in collinearity-prone panel data, cluster on \pc{} scores, not raw
      variables --- the raw space \emph{suppresses} the model-selection
      signal.
  \end{itemize}
\end{frame}


\begin{frame}{The linearity Certificate, as a Check List}
  For any small, collinearity-prone panel you intend to cluster:
  \begin{enumerate}
    \item \textbf{Measure Effective Dimensionality} --- participation ratio of
      the eigenvalue spectrum. Low \pr{} $\Rightarrow$ expect a near-linear
      manifold and a distorted raw space.
    \item \textbf{Cluster on \pc{} scores}, and compare against raw-space
      assignments (\ari{}). Explain any invariance geometrically.
    \item \textbf{Scan $k$ in the Clean Space} --- internal indices in raw
      space can be suppressed by the collinear axis.
    \item \textbf{Attack with Kernel Families that Fail Differently} ---
      distance-based, information-geometric, graph-based; median-heuristic
      bandwidths, nothing tuned. Cross-family agreement is the certificate.
    \item \textbf{Cross-check with Supervised Kernels} --- kernel vs.\ linear
      ridge under LOOCV; audit outliers for extrapolation effects before
      blaming curvature.
  \end{enumerate}
\end{frame}


\begin{frame}{Future Work}
  \begin{itemize}
    \item \textbf{Temporal Dynamics.} Hidden Markov models or regime-switching
      VAR: quantify state persistence and transitions instead of static bands.
    \item \textbf{Modern representation-learning baselines.} $t$-SNE, UMAP,
      autoencoders, spectral and contrastive methods: does the certified
      linear structure persist under representations that don't assume it?
    \item \textbf{Other transport sectors.} Public transport panels share the
      same secular-growth collinearity --- is the
      collinearity--silhouette-suppression pattern domain-general?
      \muted{(Railways have profit cycles too.)}
  \end{itemize}
\end{frame}


\begin{frame}[plain]
  \vspace*{1.1cm}
  \bigstatement{Understand your data before you model it ---\\
    and when you claim it is linear, \textcolor{ascTeal600}{certify it}.}
  \vspace{0.7cm}

  \muted{Kernel PCA, 28 years on, working as the audit tool for its own
  linear ancestor.}
  \vspace{0.8cm}

  %{\small
  %\textbf{Thanks} --- Prof.\ Renold, for discussions and for making data and
  %methodology public.\\
  %Prof.\ Bunke, for mentorship and teaching at the FKI research group, here
  %in Bern.\par}
  %\vspace{0.6cm}
\end{frame}

%% ============================================================
%% Backup slides
%% ============================================================
\appendix

\begin{frame}{Backup: Full Cluster-Quality Table}
  \begin{center}
  \setlength{\tabcolsep}{9pt}
  \begin{tabular}{@{}c rr rr rr@{}}
  \toprule
   & \multicolumn{2}{c}{\textbf{7D-raw}} &
     \multicolumn{2}{c}{\textbf{3D-\pc{}}} &
     \multicolumn{2}{c}{\textbf{4D-\pc{}}} \\
  \cmidrule(lr){2-3}\cmidrule(lr){4-5}\cmidrule(lr){6-7}
  $k$ & \sil{} & \db{} & \sil{} & \db{} & \sil{} & \db{} \\
  \midrule
  2 & 0.431 & 0.906 & 0.506 & 0.742 & 0.458 & 0.853 \\
  \textbf{3} & 0.443 & 0.648 & \textbf{0.523} & \textbf{0.507} & 0.472 & 0.601 \\
  4 & 0.429 & 0.658 & 0.481 & 0.603 & 0.466 & 0.615 \\
  5 & 0.413 & 0.679 & 0.421 & 0.603 & 0.460 & 0.569 \\
  6 & 0.448 & 0.647 & 0.448 & 0.584 & \textbf{0.490} & \textbf{0.557} \\
  \textbf{7} & \textbf{0.449} & \textbf{0.629} & 0.429 & 0.583 & 0.463 & 0.640 \\
  8 & 0.439 & 0.613 & 0.416 & 0.596 & 0.449 & 0.627 \\
  9 & 0.416 & 0.538 & 0.430 & 0.554 & 0.410 & 0.578 \\
  \midrule
  Optimal $k$ & \multicolumn{2}{c}{$k=7$} & \multicolumn{2}{c}{$k=3$} &
    \multicolumn{2}{c}{$k=6$} \\
  \bottomrule
  \end{tabular}
  \end{center} 
\end{frame}

\begin{frame}{Backup: Loadings on \pc{1}--\pc{4}}
  \begin{center}
  \setlength{\tabcolsep}{7pt}
  \begin{tabular}{@{}l *{7}{r}@{}}
  \toprule
   & \multicolumn{7}{c}{\textbf{Variable}} \\
  \cmidrule(l){2-8}
  \textbf{\pc{}} & Op.~Profit & Fuel & Wages & RPM & LF & Yield & Other Exp. \\
  \midrule
  1 (59.8\%) & $+0.33$ & $+0.39$ & $+0.28$ & $\mathbf{+0.45}$ &
    $\mathbf{+0.45}$ & $+0.34$ & $-0.38$ \\
  2 (16.7\%) & $\mathbf{+0.63}$ & $-0.30$ & $\mathbf{-0.51}$ & $+0.23$ &
    $+0.27$ & $-0.22$ & $+0.27$ \\
  3 (11.3\%) & $+0.24$ & $\mathbf{-0.45}$ & $+0.55$ & $+0.08$ &
    $-0.17$ & $\mathbf{+0.45}$ & $+0.46$ \\
  4 (7.8\%)  & $+0.16$ & $+0.17$ & $-0.50$ & $-0.37$ & $-0.18$ &
    $\mathbf{+0.73}$ & $-0.04$ \\
  \bottomrule
  \end{tabular}
  \end{center}
  \vspace{0.3cm}
  \begin{itemize}
    \item \pc{1} uniformly positive (except other expenses): the fingerprint
      of collinearity --- one shared growth factor.
    \item \pc{4}: yield per passenger $+0.73$, the COVID signature.
  \end{itemize}
\end{frame}



\begin{frame}{Backup: the $k=3$ vs.\ $k=6$ Tension Across Spaces}
  \begin{itemize}
    \setlength{\itemsep}{0.35cm}
    \item In 3D-\pc{}: both indices pick $k=3$. In 4D-\pc{}: both pick $k=6$.
      Not a contradiction --- \pc{4} is essentially a
      \teal{single-year outlier direction} (2020).
    \item Adding \pc{4} manufactures a well-separated singleton cluster
      $C_3 = \{2020\}$, which internal indices reward --- inflating $k=6$
      quality in exactly the space that isolates the outlier.
    \item Reading: three \emph{structural} regimes + one \emph{singularity};
      the six-cluster taxonomy subdivides the three regimes descriptively.
    \item Suggestive, not confirmed: $n=26$, two internal indices, no external
      ground truth, no significance test.
  \end{itemize}
\end{frame}



\begin{frame}{Backup: Data and Implementation Notes}
  \begin{itemize}
    \item Data digitized from Renold et al.'s published tables; 2008 is absent
      from the original data table, hence $n=25$ complete years for the
      regression experiments.
    \item All variables standardized before distance computations.
    \item \kmeans{}: scikit-learn, \texttt{n\_init=100},
      \texttt{max\_iter=1000}.
    \item Kernel ridge: predictors and target standardized \emph{within each
      training fold}; $\gamma$ and $\alpha$ by nested LOOCV grid search ---
      no test-fold leakage.
    \item Code available on request.
  \end{itemize}
\end{frame}



\begin{frame}{Backup: why not Fisher Discriminant Analysis?}
  \begin{itemize}
    \item FDA needs pre-labelled membership; the research question is
      structure \emph{discovery} --- using the labels we are trying to
      validate would be circular.
    \item With $n \approx 5$--$9$ per cluster, within-class covariance is
      near-singular; the computation destabilizes.
    \item Post-hoc, FDA/LDA and supervised \kpca{} are attractive: check
      whether supervised projections align with the unsupervised ones found
      here.
  \end{itemize}
\end{frame}



\begin{frame}{Backup: kernel definitions and bandwidths}
  \begin{itemize}
    \item \textbf{RBF:} $K(x,y)=\exp(-\gamma\|x-y\|^2)$,
      $\gamma$ from the median pairwise distance in 4D-\pc{} space.
    \item \textbf{Polynomial:} $K(x,y)=(x^{\top}y + c)^d$, $d \in \{2,3\}$.
    \item \textbf{Fisher--Mahalanobis:} $K_F(x,y)=x^{\top}\hat\Sigma^{-1}y$,
      $\hat\Sigma=\mathrm{diag}(\lambda_1,\dots,\lambda_4)$ (diagonal by PCA
      orthogonality) $\equiv$ PCA on whitened scores
      $z_k = x_k/\sqrt{\lambda_k}$; derived from a single-Gaussian generative
      model.
    \item \textbf{Graph diffusion:} $K=\exp(-\beta L)$ on the 5-NN graph of
      4D-\pc{} scores, $\beta$ from the median edge weight.
    \item \kpca{} operates on the $26\times26$ Gram matrix; cluster
      assignments held fixed from the 4D-\pc{} \kmeans{} solution.
  \end{itemize}
\end{frame}

\end{document}
